\documentclass[a4paper,14pt]{report} %размер бумаги устанавливаем А4, шрифт 14 пунктов
\usepackage[T2A]{fontenc}
\usepackage[english,russian]{babel} % включаем русский язык
\usepackage[utf8x]{inputenc} % работаем с форматом  utf-8
\usepackage{tocloft}% изменяем параметры оглавления
\usepackage{graphicx} % библиотека для работы с графикой
\graphicspath{ {./img/} } % путь к папке с изображениями
\usepackage{amssymb,amsfonts,amsmath,mathtext,cite,enumerate,float} %подключаем нужные пакеты расширений
\usepackage{cite} % Красивые ссылки на литературу
\usepackage{mathtext}
\usepackage{float}

\usepackage{textcomp}
%\usepackage{minted}      % (requires -shell-escape)
\usepackage{xcolor}
\usepackage{filecontents}
\usepackage{caption} % описания
\usepackage{color} %% это для отображения цвета в коде
\usepackage{listings} %% собственно, это и есть пакет listings
\linespread{1.3}
\usepackage{geometry} % Меняем поля страницы
\geometry{left=2cm}% левое поле
\geometry{right=1.5cm}% правое поле
\geometry{top=1cm}% верхнее поле
\geometry{bottom=2cm}% нижнее поле
%\usepackage{minted}      % (requires -shell-escape)
%%% Стили %%%
\bibliographystyle{BibTeX-Styles/utf8gost71u}    % Оформляем библиографию по ГОСТ 7.1 (ГОСТ Р 7.0.11-2011, 5.6.7)
\cftsetpnumwidth{3pt} % точки прям до цифр в оглавлении

\makeatletter
\renewcommand{\@biblabel}[1]{#1.} % Заменяем библиографию с квадратных скобок на точку:
\makeatother

\renewcommand{\theenumi}{\arabic{enumi}}% Меняем везде перечисления на цифра.цифра
\renewcommand{\labelenumi}{\arabic{enumi})}% Меняем везде перечисления на цифра.цифра
\renewcommand{\theenumii}{.\arabic{enumii}}% Меняем везде перечисления на цифра.цифра
\renewcommand{\labelenumii}{\arabic{enumi}.\arabic{enumii}.}% Меняем везде перечисления на цифра.цифра
\renewcommand{\theenumiii}{.\arabic{enumiii}}% Меняем везде перечисления на цифра.цифра
\renewcommand{\labelenumiii}{\arabic{enumi}.\arabic{enumii}.\arabic{enumiii}.}% Меняем везде перечисления на цифра.цифра
\renewcommand{\thesection}{\arabic{section}.}
\renewcommand{\thesubsection}{\arabic{section}.\arabic{subsection}.}
\newcommand{\tocsecindent}{\hspace{7mm}}
\newcommand*{\No}{\textnumero}
\definecolor{gray}{rgb}{0.5,0.5,0.5}
\addto\captionsrussian{\def\refname{Список литературы}}
\long\def\comment{}% многострочный комментарий

% документ
\begin{document}
\long\def\/*#1*/{}
	% Титульный лист (ГОСТ Р 7.0.11-2001, 5.1)
% титульный лист
\newpage
\begin{titlepage}
	\thispagestyle{empty}
	
	{\large
		\begin{center}
			%\thesisOrganization
			\textsl{Федеральное государственное бюджетное образовательное учреждение высшего профессионального образования \\* }	
			\vspace{2em}
		\end{center}
	}
	%\vspace{0pt plus4fill} %число перед fill = кратность относительно некоторого расстояния fill, кусками которого заполнены пустые места
	\IfFileExists{img/bmstu-logo.png}
	{
		\begin{minipage}[t]{0.13\textwidth}
			\centering\raisebox{\dimexpr \topskip-\height}{%
				\includegraphics[width=\textwidth]{bmstu-logo.png}}
			\label{fig0}
		\end{minipage}\hfill
		\begin{minipage}[t]{0.65\textwidth}
			\begin{center}
				\large  \textsl{\textbf {\\ «Московский государственный технический университет имени Н.Э. Баумана» \\* (МГТУ им. Н.Э. Баумана) \\* }}
			\end{center}
		\end{minipage}
	}
	
	\vspace{2em}
	
	\hrulefill
	
	\begin{center}
		\vspace{0pt plus2fill} %число перед fill = кратность относительно некоторого расстояния fill, кусками которого заполнены пустые места
		{\large
			Факультет  <<Информатика и системы управления>>\\
			Кафедра  <<Программное обеспечение ЭВМ и информационные технологии>>
		}
	\end{center}
	
	{\Large
		\begin{center}
			\textbf{Отчет по лабораторной работе №17 \\* по курсу: \\* } «Функциональное и Логическое программирование» \\
		\end{center}
	}
	%
	\vspace{0pt plus4fill} %число перед fill = кратность относительно некоторого расстояния fill, кусками которого заполнены пустые места
	
	\begin{center}
		%Научный руководитель:
		\newlength{\ML}
		\hfill Студент группы ИУ7-63Б:\hspace{0.2cm} Фурдик Н.\,О.\\
		\hfill\hspace{1.7cm} (Фамилия И.О.)\\
		\hfill Преподаватель:\hspace{0.2cm} Толпинская Н.\,Б., Строганов Ю.\,В.\\
		\hfill\hspace{1.7cm} (Фамилия И.О.)\\
		%\supervisorRegalia
		
		%\supervisorFio
	\end{center}
	%
	
	\vspace{0pt plus4fill} %число перед fill = кратность относительно некоторого расстояния fill, кусками которого заполнены пустые места
	{\centering Москва, 2020\par}
\end{titlepage} % титульный лист
	\newpage
	\tableofcontents % это оглавление, которое генерируется автоматически
	% --- ugly internals for language definition ---
%
\makeatletter
\lstset{
	captionpos = below,
	frame      = single,
	columns    = fullflexible,
	basicstyle = \small\sffamily,
}
% initialisation of user macros
\newcommand\PrologPredicateStyle{}
\newcommand\PrologVarStyle{}
\newcommand\PrologAnonymVarStyle{}
\newcommand\PrologAtomStyle{}
\newcommand\PrologOtherStyle{}
\newcommand\PrologCommentStyle{}

% useful switches (to keep track of context)
\newif\ifpredicate@prolog@
\newif\ifwithinparens@prolog@

% save definition of underscore for test
\lst@SaveOutputDef{`_}\underscore@prolog

% local variables
\newcount\currentchar@prolog

\newcommand\@testChar@prolog%
{%
	% if we're in processing mode...
	\ifnum\lst@mode=\lst@Pmode%
	\detectTypeAndHighlight@prolog%
	\else
	% ... or within parentheses
	\ifwithinparens@prolog@%
	\detectTypeAndHighlight@prolog%
	\fi
	\fi
	% Some housekeeping...
	\global\predicate@prolog@false%
}

% helper macros
\newcommand\detectTypeAndHighlight@prolog
{%
	% First, assume that we have an atom.
	\def\lst@thestyle{\PrologAtomStyle}%
	% Test whether we have a predicate and modify the style accordingly.
	\ifpredicate@prolog@%
	\def\lst@thestyle{\PrologPredicateStyle}%
	\else
	% Test whether we have a predicate and modify the style accordingly.
	\expandafter\splitfirstchar@prolog\expandafter{\the\lst@token}%
	% Check whether the identifier starts by an underscore.
	\expandafter\ifx\@testChar@prolog\underscore@prolog%
	% Check whether the identifier is '_' (anonymous variable)
	\ifnum\lst@length=1%
	\let\lst@thestyle\PrologAnonymVarStyle%
	\else
	\let\lst@thestyle\PrologVarStyle%
	\fi
	\else
	% Check whether the identifier starts by a capital letter.
	\currentchar@prolog=65
	\loop
	\expandafter\ifnum\expandafter`\@testChar@prolog=\currentchar@prolog%
	\let\lst@thestyle\PrologVarStyle%
	\let\iterate\relax
	\fi
	\advance \currentchar@prolog by 1
	\unless\ifnum\currentchar@prolog>90
	\repeat
	\fi
	\fi
}
\newcommand\splitfirstchar@prolog{}
\def\splitfirstchar@prolog#1{\@splitfirstchar@prolog#1\relax}
\newcommand\@splitfirstchar@prolog{}
\def\@splitfirstchar@prolog#1#2\relax{\def\@testChar@prolog{#1}}

% helper macro for () delimiters
\def\beginlstdelim#1#2%
{%
	\def\endlstdelim{\PrologOtherStyle #2\egroup}%
	{\PrologOtherStyle #1}%
	\global\predicate@prolog@false%
	\withinparens@prolog@true%
	\bgroup\aftergroup\endlstdelim%
}

% language name
\newcommand\lang@prolog{Prolog-pretty}
% ``normalised'' language name
\expandafter\lst@NormedDef\expandafter\normlang@prolog%
\expandafter{\lang@prolog}

% language definition
\expandafter\expandafter\expandafter\lstdefinelanguage\expandafter%
{\lang@prolog}
{%
	language            = Prolog,
	keywords            = {},      % reset all preset keywords
	showstringspaces    = false,
	alsoletter          = (,
	alsoother           = @$,
	moredelim           = **[is][\beginlstdelim{(}{)}]{(}{)},
	MoreSelectCharTable =
	\lst@DefSaveDef{`(}\opparen@prolog{\global\predicate@prolog@true\opparen@prolog},
}

% Hooking into listings to test each ``identifier''
\newcommand\@ddedToOutput@prolog\relax
\lst@AddToHook{Output}{\@ddedToOutput@prolog}

\lst@AddToHook{PreInit}
{%
	\ifx\lst@language\normlang@prolog%
	\let\@ddedToOutput@prolog\@testChar@prolog%
	\fi
}

\lst@AddToHook{DeInit}{\renewcommand\@ddedToOutput@prolog{}}

\makeatother
%
% --- end of ugly internals ---


% --- definition of a custom style similar to that of Pygments ---
% custom colors
\definecolor{PrologPredicate}{RGB}{000,031,255}
\definecolor{PrologVar}      {RGB}{024,021,125}
\definecolor{PrologAnonymVar}{RGB}{000,127,000}
\definecolor{PrologAtom}     {RGB}{186,032,032}
\definecolor{PrologComment}  {RGB}{063,128,127}
\definecolor{PrologOther}    {RGB}{000,000,000}

% redefinition of user macros for Prolog style
\renewcommand\PrologPredicateStyle{\color{PrologPredicate}}
\renewcommand\PrologVarStyle{\color{PrologVar}}
\renewcommand\PrologAnonymVarStyle{\color{PrologAnonymVar}}
\renewcommand\PrologAtomStyle{\color{PrologAtom}}
\renewcommand\PrologCommentStyle{\itshape\color{PrologComment}}
\renewcommand\PrologOtherStyle{\color{PrologOther}}

% custom style definition 
\lstdefinestyle{Prolog-pygsty}
{
	language     = Prolog-pretty,
	upquote      = true,
	stringstyle  = \PrologAtomStyle,
	commentstyle = \PrologCommentStyle,
	literate     =
	{:-}{{\PrologOtherStyle :-}}2
	{,}{{\PrologOtherStyle ,}}1
	{.}{{\PrologOtherStyle .}}1
}



\newpage
\section*{Постановка задачи}
\addcontentsline{toc}{section}{\tocsecindent{Постановка задачи}}

\Large В одной программе написать правила, позволяющие найти 
\begin{enumerate}
	\item \textbf{Максимум из двух чисел} 	
	\begin{itemize}  \item без использования отсечения;
		\item с использованием отсечения.\end{itemize}
	\item \textbf{Максимум из трех чисел}  	\begin{itemize}  \item без использования отсечения;
		\item с использованием отсечения.\end{itemize}
\end{enumerate}
Убедиться в правильности результатов.
Для каждого случая пункта 2 обосновать необходимость всех условий тела. 
Для одного из вариантов ВОПРОСА и каждого варианта задания 2 составить таблицу, отражающую конкретный порядок работы системы: \\
Т.к. резольвента хранится в виде стека, то состояние резольвенты требуется отображать в столбик: вершина – сверху! Новый шаг надо начинать с нового состояния резольвенты!


\newpage
\section*{Листинг программы}
\addcontentsline{toc}{section}{\tocsecindent{Листинг программы}}
Ниже представлен листинг программы:
\lstinputlisting[
style      = Prolog-pygsty,
caption    = {Задания 1 и 2},
]{src/1.pro}

\section*{Описание работы системы}
\addcontentsline{toc}{section}{\tocsecindent{Описание порядка поиска ответов}}

Ниже представлен алгоритм поиска ответов на вопросы max\_of\_three(Max, 1, 3, 7) и max\_of\_three(Max, 3, 7, 1) с использованием отсечения и без него.

\begin{table}
	\caption{Описание работы системы без использования отсечения}
	\begin{tabular}{|p{1cm}|p{5cm}|p{5cm}|p{5cm}|}
		\hline
		№ шага & Состояние резольвенты, и вывод: дальнейшие действия (почему?) & Для каких термов запускается алгоритм унификации: Т1=Т2 и каков результат (и подстановка) & Дальнейшие действия: прямой ход или откат (почему и к чему приводит?) \\
		\hline
		1 & Резольвента: max\_of\_three(Max, 1, 3, 7). Начинается поиск совпадений по БЗ &  & прямой ход \\
		\hline
		2 & Резольвента: max\_of\_three(Max, 1, 3, 7).  & Нашли подходящее правило: max\_of\_three(A, B, C, D) :- C > D, C >= B, A = C. Подставляем A = Max, B = 1, C = 3, D = 7 & прямой ход \\
		\hline
		3 & Резольвента: C > D,\newline max\_of\_three(Max, 1, 3, 7). & Пробуем унифицировать: 3 > 7, не подходит. Откат. & откат, продолжаем поиск \\
		\hline
		4 & Резольвента: max\_of\_three(Max, 1, 3, 7). & Нашли подходящее правило: max\_of\_three(A, B, C, D) :- B >= D, B > C, A = B.. Подставляем A = Max, B = 1, C = 3, D = 7 & прямой ход \\
		\hline
		5 & Резольвента: B >= D, \newline max\_of\_three(Max, 1, 3, 7). & Пробуем унифицировать: 1 >= 7, не подходит. Откат. & откат, продолжаем поиск \\
		\hline
		6 & Резольвента: max\_of\_three(Max, 1, 3, 7). & Нашли подходящее правило: max\_of\_three(A, B, C, D) :- D >= B, D >= C, A = D. Подставляем A = Max, B = 1, C = 3, D = 7 & прямой ход \\
		\hline
		7 & Резольвента: D >= B, \newline max\_of\_three(Max, 1, 3, 7). & Пробуем унифицировать: 7 >= 1, подходит, идем дальше.   & прямой ход \\
		\hline
		8 & Резольвента: D >= C,\newline max\_of\_three(Max, 1, 3, 7). & Пробуем унифицировать: 7 >= 3, подходит, идем дальше. & прямой ход \\
		\hline
		9 & Резольвента: A = D, \newline max\_of\_three(Max, 1, 3, 7). & Связываем A = D = 7.  & прямой ход \\
		\hline
		10 & max\_of\_three(Max, 1, 3, 7) & Подставляем Max = 7. Ответ найден & прямой ход \\
		\hline
		Вывод: & \textbf{Max = 7} &  &  \\
		\hline
	\end{tabular}
\end{table}

\begin{table}
	\caption{Описание работы системы c использованием отсечения}
	\begin{tabular}{|p{1cm}|p{5cm}|p{5cm}|p{5cm}|}
		\hline
		№ шага & Состояние резольвенты, и вывод: дальнейшие действия (почему?) & Для каких термов запускается алгоритм унификации: Т1=Т2 и каков результат (и подстановка) & Дальнейшие действия: прямой ход или откат (почему и к чему приводит?) \\
		\hline
		1 & Резольвента: max\_of\_three(Max, 3, 7, 1). Начинается поиск совпадений по БЗ &  & прямой ход \\
		\hline
		2 & Резольвента: max\_of\_three(Max, 3, 7, 1). & Нашли подходящее правило: max\_of\_three(A, B, C, D) :- C >= D, C >= B, A = C, !. Подставляем A = Max, B = 1, C = 3, D = 7 & прямой ход \\
		\hline
		3 & Резольвента: C >= D, \newline max\_of\_three(Max, 3, 7, 1). & Пробуем унифицировать: 7 >= 1, подходит, идем дальше.   & прямой ход \\
		\hline
		4 & Резольвента: C >= B,\newline max\_of\_three(Max, 3, 7, 1). & Пробуем унифицировать: 7 >= 3, подходит, идем дальше. & прямой ход \\
		\hline
		5 & Резольвента: A = D, \newline max\_of\_three(Max, 3, 7, 1). & Связываем A = D = 7.  & прямой ход \\
		\hline
		6 & max\_of\_three(Max, 3, 7, 1) & Связываем Max = A = 7. Ответ найден. Поскольку отсечение есть, дальше не идем. & прямой ход \\
		\hline
		Вывод: & \textbf{Max = 7} &  &  \\
		\hline
	\end{tabular}
\end{table}
	% ответы на вопросы
\newpage
\section*{Ответы на вопросы}
\addcontentsline{toc}{section}{\tocsecindent{Ответы на вопросы}}
\begin{enumerate}
    \item \textbf{В какой части правила сформулировано знание? Это знание о чем, с формальной точки зрения?}\\
    Знание формулируется в правой части правила. Оно содержит в себе некую информацию, которая является истиной.
    
    \item \textbf{Что такое процедура?}\\
    Процедурой называется совокупность правил, заголовки которых имеют одно и то же имя и одну и ту же арность (местность), т.е. это совокупность правил, описывающих одно определенное отношение. Процедура описывает одно отношение, которое нельзя сформулировать одним предложением.
    
    \item \textbf{Сколько в БЗ  текущего задания процедур?}\\
    В БЗ текущего задания четыре процедуры.

    \item \textbf{Что такое пример терма, это частный случай терма, пример? Как строится пример?}\\
    Терм B называется примером терма $A$, если существует такая подстановка такая, что $B = A  \theta$.
    
    Терм C является общим примером $A, B$ если существует подстановки $\theta_1$ и $\theta_2$ такие что С$ С = A \theta_1$ и $C = B \theta_2$.
	
	Пример. Имеется БЗ, в которой есть факт g(1, a, b). Базе Знаний подаётся запрос g(x, y, z). Cистема находит общий пример и только тогда утверждение из Базы Знаний может считаться доказательством истинности.
	
	
    \item \textbf{Что такое наиболее общий пример?}\\
   	Терм $S$ называется более общим чем терм $Т$, если $T$ является примером $S$, а $S$ не является примером $T$. 
   	$S$ называется наиболее общим примером $T_1$ и $T_2$, если $S$ такой их общий пример, который является более общим по отношению к любому другому их примеру. 
    
    \item \textbf{Назначение и результат работы алгоритма унификации. Что значит двунаправленная передача параметров при работе алгоритма унификации, поясните на примере одного из случаев пункта  3.}\\
    Алгоритм встроен в систему и для своей работы использует три области памяти:
    \begin{itemize}
    	
    
    \item стек: хранит равенства, унификацию которых надо доказать;
    \item рабочее поле:  равенство вытащенное из стека, которое сейчас доказываем;
    \item результирующая ячейка памяти: накапливает наиболее общий унификатор;
    \item флаг - переменная 'неудача':  равна 1, если возникла тупиковая ситуация иначе 0.
	\end{itemize}
    \textbf{шаг:}
    \textbf{начало:} Сопоставляем $TТ_1$=$ТT_2$\\
    занести в стек $TТ_1$=$ТT_2$\\
    положить неудача=0\\
    пока стек не пуст – \textbf{цикл}:
    \begin{itemize}
    \item	считать из стека в рабочую область очередное равенство S=Т
    \item	обработать считанное по правилам:
    \begin{itemize}
    \item	если S и Т несовпадающие константы,			то неудача=1, и выход из цикла
    \item	если одинаковые константы					то следующий шаг цикла
    \item	если S переменная и Т терм содержащий S,		то неудача=1, и выход из цикла
    \item	если S переменная и Т терм НЕ содержащий S,	то отыскать в стеке и в результирующей ячейке все вхождения S и заменить на Т. Добавить в результирующую ячейку равенство S=Т.  Следующий шаг цикла
    \item	если S и Т составные термы с разными функторами или разными арностями, то неудача=1, выход из цикла
    \item	если S и Т составные термы с одинаковыми функторами и арностью: $S=f(s_1 s_2 .. s_m);$   $T=f(t_1 t_2 ... t_m),$  то занести в стек равенство $S_1=T_1, S_2=T_2 ... S_m=T_m.$
\end{itemize}
    \item	очистить рабочее поле
	\end{itemize}
    –  \textbf{конец цикла}
    
    \textbf{если} неудача = 1 то унификация невозможна
    \\ \textbf{если} неудача = 0 то унификация успешна, результирующая ячейка содержит наиболее общий унификатор
    \\ \textbf{конец.}
    
    Таким образом, с помощью алгоритма унификации происходит двунаправленная передача параметров процедурам. Например, из внешнего мира в программу для дальнейшего использования или из программы во внешний мир – значения интересующего нас параметра. 
    
    
    
    \item \textbf{В каком случае запускается механизм отката?}\\
    Для обоснования истинности вопроса, не всегда бывает возможным подобрать нужное знание в принципе, или не всегда это можно выполнить однозначно. Кроме этого, не очевиден принцип и алгоритм (последовательность перебора) выбора знания для использования. Часто при выборе знания  и человеку и системе необходимо оценить содержательный смысл – семантику утверждения, что на формальном уровне однозначно выполнить сложно. Следует заметить, что недетерминизм поиска ответа – это базовый принцип работы естественного, искусственного интеллекта и системы Prolog. Таким образом, в сложных ситуациях – при высокой степени неопределенности, и человек и система Prolog совершают ошибки. Для того, что бы такая ошибка выбора знания для системы Prolog была не критична, в систему встроен механизм отката (на момент последнего сделанного выбора) и отмена принятого решения (выполненной конкретизации).
    
    \item \textbf{Виды и назначение переменных в Prolog. Примеры из задания.  Почему использованы те или другие переменные (примеры из задания)?}\\
    Переменные начинаются с заглавной буквы или с символа подчёркивания и могут быть анонимными или именованными (анонимная переменная обозначается нижним подчеркивание). Говорят, что переменная может быть связана с некоторым значением или оставаться независимой. Именованные переменные уникальны в рамках одного предложения. Анонимная переменная уникальна всегда. Переменные предназначены для передачи значений «во времени и в пространстве».

\end{enumerate}

% ответ
%\\
	% список литературы
\newpage
\addcontentsline{toc}{section}{\tocsecindent{Список литературы}}

\begin{thebibliography}{}
    \bibitem{litlink1}  Население Земли. [Электронный ресурс]. – Режим доступа: http://www.countrymeters.info (дата обращения 30.05.2020)
    \bibitem{litlink2}  Дефицит пресной воды в странах мира. Справка. [Электронный ресурс]. – Режим доступа:  https://ria.ru/20100322/215718166.html (дата обращения 30.05.2020)
     \bibitem{litlink3} ДЕФИЦИТ ПРЕСНОЙ ВОДЫ: ПРОБЛЕМЫ И СПОСОБЫ РЕШЕНИЯ [Электронный ресурс]. – Режим доступа: 
    https://thewallmagazine.ru/lack-of-fresh-water/ (дата обращения 30.05.2020)
  \bibitem{litlink4} Water scarcity [Электронный ресурс]. – Режим доступа:  https://en.wikipedia.org/wiki/Water\_scarcity (дата обращения 30.05.2020)
  \bibitem{litlink5}Конференция "Международное десятилетие снабжения питьевой водой и санитарии". [Электронный ресурс]. – Режим доступа:  https://documents-dds-ny.un.org/doc/RESOLUTION/GEN/NR0/394/93/IMG/NR039493.pdf?OpenElement 
  \bibitem{litlink6} Всемирный доклад ООН.[Электронный ресурс]. – Режим доступа: \\ https://unesdoc.unesco.org/ark:/48223/pf0000231823\_eng
\end{thebibliography}

\end{document}